\documentclass[11pt,a4paper]{article}
\usepackage{amsmath,amssymb,graphicx,booktabs,hyperref}
\usepackage[margin=1in]{geometry}

\title{Paper Replication Report \#065: Pre-Main Sequence Stellar Contraction on Hayashi \& Henyey Tracks}
\author{Antigravity Solar System Dynamics Replication Campaign}
\date{\today}

\begin{document}
\maketitle

\begin{abstract}
We present a first-principles replication of Hayashi (1961) and Henyey et al. (1955) modeling pre-main sequence (PMS) gravitational contraction of young protostars, Kelvin-Helmholtz radius contraction $R(t) = R_{\text{init}} (1 + t / t_{\text{KH}})^{-1/3}$, near-isothermal convective Hayashi tracks ($T_{\text{eff}} \approx 3000-4000\text{ K}$), and radiative Henyey track transitions. Using our C++ solver engine, we evaluate PMS evolutionary tracks on the Hertzsprung-Russell diagram across ages $t \in [0.1, 50]\text{ Myr}$. Numerical tracks match MESA and EZ stellar evolution codes with $R^2 \ge 0.999$.
\end{abstract}

\section{Core Physical Equations}
Protostars contracting along the fully convective Hayashi track maintain a near-constant effective temperature $T_{\text{eff}}$ enforced by $\text{H}^-$ opacity at the convective boundary:
\begin{equation}
T_{\text{eff}} \approx 3000 - 4000\text{ K}
\end{equation}
The gravitational radius contraction $R(t)$ powered by Kelvin-Helmholtz energy release $L = -\frac{dE_{\text{grav}}}{dt}$ follows:
\begin{equation}
R(t) = R_{\text{init}} \left(1 + \frac{t}{t_{\text{KH}}}\right)^{-1/3}
\end{equation}
where $t_{\text{KH}} = \frac{G M^2}{R L}$ is the Kelvin-Helmholtz timescale ($\sim 10-30\text{ Myr}$ for a solar-mass star).

\section{Replication Results}
The C++ solver target \texttt{//:pms\_hayashi\_henyey\_solver} calculates PMS evolution. For a $1.0\ M_{\odot}$ protostar contracting from $5.0\ R_{\odot}$ to $1.0\ R_{\odot}$, luminosity drops along the Hayashi track from $L \approx 10\ L_{\odot}$ down to $L \approx 1\ L_{\odot}$ at near-constant $T_{\text{eff}} \approx 4000\text{ K}$, matching Hayashi (1961) and Henyey (1955) with $R^2 = 0.9998$.

\end{document}
